\documentclass[11pt]{article}

\usepackage[margin=2.2cm]{geometry}
\usepackage{graphicx}
\usepackage{amsmath}
\usepackage{amssymb}
\usepackage{float}
\usepackage{caption}
\usepackage{enumitem}
\setlist{itemsep=2pt, topsep=4pt}

\title{QML on TR1 South: summary}
\author{Weak Lensing QML pipeline}
\date{\today}

\begin{document}
\maketitle

\section*{What we did}
\begin{itemize}
\item Applied the QML power-spectrum estimator (Maraio, Hall \& Taylor 2022) to the real \textbf{TR1 South} shear
catalogue, at $N_{\rm side}=64$.
\item \textbf{All six tomographic bins combined} into one sample (21.9 galaxies/arcmin$^2$), because single bins
(3--4 galaxies/arcmin$^2$) are much noisier.
\item Compared QML with two other estimators on the same data: \textbf{NaMaster} (pseudo-$C_\ell$) and
\textbf{heracles}, run directly on the catalogue.
\item The patch is small, 0.9\% of the sky (about $20^\circ$ across), so the spectra are measured in bands of
$\Delta\ell=20$.
\end{itemize}

\section*{1. Pixels removed}
\begin{itemize}
\item Many pixels on the edge of the survey, and along gaps inside it, are only partly covered. They contain few
galaxies, so they are much noisier than the rest: up to $\sim$13{,}000 times the average.
\item QML assumes \textbf{every pixel has the same noise}, so these pixels ruined the QML result.
\item We keep only pixels that are \textbf{at least 50\% covered}. This removes 96 of 538 pixels, but only 5\% of the
galaxies.
\item With the cut, QML works. heracles gives almost the same answer with or without the cut, because it already
down-weights partly covered pixels.
\end{itemize}

\begin{figure}[H]
\centering
\includegraphics[width=0.8\textwidth]{summary_figures/pixels.png}
\caption{Left: pixels kept (blue) and removed (orange). Right: how much of each pixel is covered; the line is the
50\% cut.}
\end{figure}

\section*{2. E-modes: all three estimators agree}
QML, NaMaster and heracles all follow the expected (fiducial) spectrum.

\begin{figure}[H]
\centering
\includegraphics[width=0.75\textwidth]{summary_figures/ee.png}
\caption{TR1 South $E$-mode power spectrum from QML, NaMaster and heracles, with the fiducial model.}
\end{figure}

\section*{3. QML has smaller error bars}
\begin{itemize}
\item The plot shows NaMaster's error bars divided by QML's. Above 1, QML is more precise.
\item On large scales ($\ell\approx20$), QML's errors are \textbf{about 3 times smaller for E-modes} and
\textbf{about 28 times smaller for B-modes}.
\item The two become similar by $\ell\approx100$, as the QML paper predicts: QML's advantage is on large scales.
\end{itemize}

\begin{figure}[H]
\centering
\includegraphics[width=0.7\textwidth]{summary_figures/sigma.png}
\caption{Ratio of NaMaster to QML error bars, for E-modes (blue) and B-modes (green).}
\end{figure}

\section*{4. B-modes}
\begin{itemize}
\item Weak lensing produces essentially no B-modes, so the measured B-modes should just be noise. Anything above the
noise points to leakage or a systematic.
\item On large scales ($\ell\lesssim50$), \textbf{QML stays close to the noise level}, while NaMaster and heracles are
about 10 times higher. That excess is E-mode power leaking into B through the small, irregular mask, which QML avoids.
\item At $\ell\gtrsim110$ all three estimators are somewhat above the noise. Because all three see it, it is in the
data or the pixelisation, not in QML. It is still an open question.
\end{itemize}

\begin{figure}[H]
\centering
\includegraphics[width=0.75\textwidth]{summary_figures/bb.png}
\caption{TR1 South $B$-mode power spectrum from QML, NaMaster and heracles, with the noise level.}
\end{figure}

\section*{Next steps}
\begin{itemize}
\item \textbf{Variable noise in QML:} give QML the noise of each pixel instead of one average value. Then no pixels
need removing, and QML handles the coverage the way heracles already does.
\item \textbf{Per-bin analysis:} repeat for each tomographic bin separately.
\end{itemize}

\end{document}
